\documentclass[12pt]{article}
\begin{document}

\title{Interpret Regression Coefficient Estimates}
\author{Robert Gulotty}

\maketitle

\subsection{Level-level Regression}

$$y=\beta_0+\beta_1x+\epsilon$$

$$\frac{\partial y}{\partial x}=\beta_1$$

This is to say, if x rises by one unit, y increases by one unit.


\subsection{Log-level Regression}

$$ln(y)=\beta_0+\beta_1x+\epsilon$$
Using the chain rule on the left hand side:
$$\frac{\partial ln(y)}{\partial x}=\frac{1}{y}*\frac{\partial y}{\partial x}=\frac{\frac{\partial y}{\partial x}}{y}$$
This means that $\frac{\partial ln(y)}{\partial x}$ is the change in y divided by the level of y.  If we multiply by 100, we get the percentage change.

$$\frac{\partial \beta_0+\beta_1x+\epsilon}{\partial x}=\beta_1$$

$$\frac{\frac{\partial y}{\partial x}}{y}=\beta_1$$
$$\frac{\frac{\partial y}{\partial x}}{y}*100=\beta_1*100$$

So if you take $\beta_1$, multiply it by 100, a one unit increase in x gives you a $\beta_1*100$ percentage change in y.



\subsection{Level-log Regression}

$$y=\beta_0+\beta_1ln(x)+\epsilon$$

If we increase x by one percent we expect y to increase by $\beta_1/100$ units of y.

\subsection{Level-level Regression}

$$ln(y)=\beta_0+\beta_1ln(x)+\epsilon$$



\end{document}
